\documentclass{ltjsarticle}
\usepackage{multicol}
\setlength{\columnseprule}{0.4pt}
\usepackage{graphicx}
\title{多段組の文書}
\author{山田 太郎}
\date{\today}
\begin{document}
\maketitle
\section{多段組の文書}
まずは2段組から。
\begin{multicols}{2}
ここは2段組です。こんなふうに、文書中の好きなところだけ段組にできるので便利です。こうやってつらつらと文章を書いていると、自動的に2段組の部分の行数が調整され、2つの段に分かれて文章が表示されていますね。
\end{multicols}
さぁ、次は3段組です！
\begin{multicols}{3}
ということで次は3段組です。3段組くらいまでなら、ギリギリ必要になることもあるかなぁと思います。段組は多すぎると結局文章全体が見づらくなってしまうので、何事も程々に、という感じですね。ほら、ちゃんと3段組になってるでしょう？
\end{multicols}
いよいよ4段組です！
\begin{multicols}{4}
4段組スタート！しかし、これは少なくとも僕は使ったことがないなぁというのが正直なところ。1段に収まる文章の量も少ないですし、お世辞にも「見やすい！」とは言えないですね。まぁ、一応こういう段組もサポートはされていますよということだけ覚えておけばいいのではないでしょうか。では、おつかれさまでした。
\end{multicols}
\end{document}

